% Section .............................................................................................................
\section{Two Layers in Every Physical Model}
\label{section:layers}

The collapse calculation exemplifies a general structure present in all physical models, but rarely made explicit. Every model carries at least two distinguishable layers of content:

\begin{enumerate}
  \item \textbf{A predictive-mathematical layer:} the equations, boundary conditions, and inference rules that generate quantitative predictions about measurable quantities.
  \item \textbf{An ontological layer:} the implicit commitments about what kinds of entities exist, what properties they possess independently of measurement, and how they relate causally.
\end{enumerate}

Standard empiricist philosophy of science focuses almost exclusively on the first layer \cite{vanfraassen_1980_image}. A model is adequate if its predictions agree with experiment; inadequate otherwise. On this view, the classical atomic model fails because it predicts a collapse timescale of $\sim10^{-11}$\,s while atoms are stable indefinitely. That diagnosis is correct, but incomplete.

% Subsection ..........................................................................................................
\subsection{The Ontological Layer of the Classical Model}
\label{subsection:layers}

The deeper diagnosis is that the classical model fails at the second layer. Newton's equations for dynamics, Coulomb's law for electromagnetism, and Larmor's radiation formula are not wrong in themselves---they remain excellent approximations for macroscopic charged objects. Yet, they fail for the electron in the hydrogen atom because they embed a specific ontology:

\begin{itemize}
  \item Electrons are \emph{point particles} with definite position and
    momentum at every instant.
  \item They move along \emph{continuous trajectories} governed by
    deterministic force laws.
  \item They act as \emph{classical sources} for the electromagnetic
    field, coupling through their charge and acceleration.
\end{itemize}

It is known nowadays that each of these assumptions is false. A model can therefore be \emph{mathematically consistent}, \emph{empirically adequate in its domain}, and \emph{ontologically erroneous}. The classical hydrogen atom is a case study in exactly this combination.



% Subsection ..........................................................................................................
\subsection{The General Diagnostic Distinction}
\label{subsection:distinction}

We propose the following general principle:

\begin{quote}
  \emph{When a physical model fails empirically at the boundary of its domain, one must ask whether the failure is predictive or ontological. A predictive failure can often be corrected by modifying equations within the same ontological framework. An ontological failure requires replacing the picture of what exists.}
\end{quote}

As an example, the transition from Newtonian to relativistic mechanics is largely \emph{predictive}: the ontology of point particles with definite trajectories is preserved; only force laws and the metric of spacetime change. The transition from classical to quantum mechanics is \emph{ontological}: the very category of ``particle with a definite trajectory'' ceases to apply.

\medskip

This distinction has implications for how we interpret theoretical progress. Kuhnian paradigm shifts \cite{kuhn_1962_structure} and Lakatosian research programmes \cite{lakatos_1978_methodology} both capture aspects of it, but neither systematically separates the predictive from the ontological layer. The collapse calculation makes the separation vivid: the predictive machinery is flawless; it is the ontological presuppositions that must go.

\medskip

\TODO{Expand: Cartwright on laws vs.\ capacities; Giere on model-world relations; Chakravartty on semi-realism. Key point: none of these frameworks isolates the ontological layer as a separate diagnostic target in quite the way we propose.}
